\section{Reproducción de experimentos y publicación}

\subsection{Del pipeline de datos a un modelo desplegable}

Reproducir el experimento implica tres etapas: recolección de datos, ciclo de entrenamiento, guard de despliegue. En cada una de ellas hay que monitorear continuamente las metrics (KL, entropy, critic loss) y mantener un registro de experimentos para comparar los distintos run.

\begin{itemize}
    \item \textbf{Datos}: usar un deterministic rollout para guardar state/action/reward/truncation;
    \item \textbf{Entrenamiento}: construir el actor/critic/stage scheduler, con guardado automático del best checkpoint;
    \item \textbf{Despliegue}: usar el conservative critic como guard en inference-time, manteniendo una fallback policy.
\end{itemize}

\begin{importantbox}{Los tres pasos del pipeline de reproducción}
1) primero asegurarse de que el replay buffer pueda serializarse; 2) mantener el actor update y el critic update en el mismo config; 3) cuando el guard stage dispare un rollback, registrar el KL y la entropy de ese momento.
\end{importantbox}

\subsection{Gestión de versiones y estrategia de rollback}

Para poder hacer rollback con rapidez, el curso recomienda usar git tag para registrar el config (learning rate, lambda, entropy bonus) correspondiente a cada checkpoint. Si un new run presenta un KL alto, se hace revert al tag, para evitar un merge erróneo.

\begin{warningbox}{Errores comunes en la estrategia de rollback}
No hay que reiniciar el replay buffer durante el rollback, pues se perdería el critic signal anterior; basta con volver al checkpoint, conservar el estado del buffer y volver a hacer warm up de la entropy.
\end{warningbox}

\subsection{Resumen de la sección}

Reproducir el experimento consiste en empaquetar el entrenamiento, el registro y el rollback en un pipeline ejecutable, de manera que varios equipos puedan hacer copy/paste del mismo buen resultado.

